\chapter{Penis Irritation}
\index{Penis Irritation}

\medskip
\noindent\textit{This chapter is for general education. Irritation has many possible causes and should not be diagnosed from appearance alone.}

\section*{Possible causes}
Irritation may result from friction, soap or other products, balanitis, thrush, an STI, eczema or psoriasis, a tight foreskin, diabetes, or less common skin conditions. Symptoms may include redness, itching, soreness, swelling, cracking, discharge, unpleasant smell, or pain when urinating.

\section*{What to do}
Stop suspected irritants and wash gently with warm water; dry without rubbing. If the foreskin retracts comfortably, clean underneath and replace it. Avoid sex until symptoms are assessed if an STI is possible. Do not apply steroid, antifungal, antibiotic, or anaesthetic creams unless a clinician or pharmacist recommends them.

\section*{Assessment and treatment}
A clinician may examine the skin, ask about sexual and medical history, test for infection, and check for diabetes or another cause. Treatment depends on the cause and may include a prescribed cream, antibiotics, antifungals, or treatment for a skin disorder.

\section*{When to Seek Medical Care}
Seek emergency help if you cannot pass urine, the foreskin is trapped behind the glans, swelling is rapidly increasing, or there is severe pain, fever, or an expanding infection. Arrange review for persistent or recurrent irritation, discharge, sores, bleeding, or a new lump.

\section*{Sources and further reading}
\begin{itemize}
\item NHS. \textit{Balanitis}. \url{https://www.nhs.uk/conditions/balanitis/}
\end{itemize}
